\documentclass[12pt,a4paper]{article}
\usepackage{iftex}
\usepackage[fontset=none]{ctex}
\ifXeTeX
  \setCJKmainfont{Noto Serif CJK SC}
\else
  \ifLuaTeX
    \setCJKmainfont{Noto Serif CJK SC}
  \fi
\fi

\usepackage[margin=1in]{geometry}
\usepackage{amsmath,amssymb,bm}
\usepackage{graphicx}
\usepackage{hyperref}
\usepackage{fancyhdr}
\usepackage{enumitem}
\usepackage{booktabs}
\usepackage{algorithm}
\usepackage{algpseudocode}

\hypersetup{colorlinks,linkcolor=blue,citecolor=blue,urlcolor=blue}
\setlength{\headheight}{15pt}

\pagestyle{fancy}
\fancyhf{}
\fancyhead[L]{二维浅水方程的集合卡尔曼滤波与集合平方根滤波}
\fancyhead[C]{}
\fancyhead[R]{\thepage}

\newcommand{\U}{\bm{U}}
\newcommand{\F}{\bm{F}}
\newcommand{\G}{\bm{G}}
\newcommand{\q}{\bm{q}}
\newcommand{\x}{\bm{x}}
\newcommand{\y}{\bm{y}}
\newcommand{\w}{\bm{w}}
\newcommand{\I}{\bm{I}}
\newcommand{\T}{^{\mathsf T}}
\newcommand{\R}{\mathbb{R}}
\newcommand{\calH}{\mathcal{H}}
\newcommand{\calM}{\mathcal{M}}
\newcommand{\calN}{\mathcal{N}}

\title{二维浅水方程的集合卡尔曼滤波与集合平方根滤波}
\author{杨鹏亮}
\date{\today}

\begin{document}

\maketitle

\begin{abstract}
本文系统阐述应用于二维浅水方程的集合卡尔曼滤波（Ensemble Kalman Filter, EnKF）
与集合平方根滤波（Ensemble Square Root Filter, EnSRF）方法的理论基础与算法实现。
前向模式采用守恒型浅水方程的有限体积离散，空间通量为 Rusanov 格式，时间推进为
二阶 Runge--Kutta 格式。在集合数据同化框架下，本文从线性高斯 Kalman 滤波出发，
推导了随机 EnKF 的样本协方差近似与扰动观测分析公式，进而引入确定性 EnKF 的
均值--扰动分离思想和集合子空间平方根变换。对 EnSRF 部分，本文详细给出了对称
平方根构造的一般公式，并讨论了串行标量观测 EnSRF、集合变换 Kalman 滤波（ETKF）
和集合调整 Kalman 滤波（EAKF）等变体之间的数学联系。最后，通过双系统合成实验
对两种滤波器的分析精度、集合离散度和误差空间分布进行了定量比较。
\end{abstract}

\tableofcontents

%=====================================================================
\section{引言}
%=====================================================================

数据同化方法将数值模式的预报与观测信息进行最优融合，是地球物理流体动力学中
状态估计的核心工具。对于高维非线性系统，Kalman 滤波及其扩展形式需要维护和
传播庞大的状态协方差矩阵，计算代价极高。Evensen~(1994) 提出的集合卡尔曼滤波
通过有限集合成员的样本统计量来近似协方差矩阵，避免了显式传播协方差的巨额开销，
因而在海洋、大气和水文等领域得到了广泛应用 \cite{evensen1994sequential}。

在集合滤波的框架下，存在两类主要的分析更新策略。第一类是随机 EnKF
\cite{burgers1998analysis}，通过对观测添加随机扰动来保证分析集合协方差在期望
意义上满足 Kalman 分析公式。第二类是确定性集合滤波，统称为集合平方根滤波
\cite{tippett2003ensemble}，它直接对集合扰动施加确定性的平方根变换，从而避免
了扰动观测引入的额外采样噪声。常见的确定性方案包括串行 EnSRF
\cite{whitaker2002ensemble}、集合调整 Kalman 滤波（EAKF）
\cite{anderson2001ensemble} 和集合变换 Kalman 滤波（ETKF）
\cite{bishop2001adaptive} 等。

二维浅水方程是描述旋转分层流体中重力波和涡旋动力学的最简模型，也是检验数据
同化方法的理想平台。本文以二维浅水方程为前向模式，系统阐述随机 EnKF 和
确定性 EnSRF 的数学基础与数值实现。与侧重代码实现细节的工作文档不同，本文
旨在提供一份自包含的理论推导，使读者能够独立理解两类方法的数学结构、算法流程
及其相互联系。各节安排如下：第~\ref{sec:forward} 节介绍前向模式及其数值离散；
第~\ref{sec:framework} 节建立集合数据同化的一般框架；第~\ref{sec:enkf}
节和第~\ref{sec:ensrf} 节分别详述随机 EnKF 和确定性 EnSRF；第~\ref{sec:experiments}
节给出数值实验结果；第~\ref{sec:conclusion} 节总结全文；附录提供了详细的
分步算法描述。

%=====================================================================
\section{前向模式：二维浅水方程的数值离散}
\label{sec:forward}
%=====================================================================

\subsection{控制方程}

考虑平坦底地形条件下的二维浅水方程。记水深为 $h(x,y,t)$，$x$ 和 $y$ 方向的
垂向平均速度分量分别为 $u$ 和 $v$。引入守恒变量
\begin{equation}
  \U =
  \begin{pmatrix}
    h \\ hu \\ hv
  \end{pmatrix}
  =
  \begin{pmatrix}
    U_1 \\ U_2 \\ U_3
  \end{pmatrix},
  \label{eq:state}
\end{equation}
其中 $hu$ 和 $hv$ 为两个方向的动量。守恒形式的浅水方程为
\begin{equation}
  \partial_t \U + \partial_x \F(\U) + \partial_y \G(\U) = 0,
  \label{eq:swe}
\end{equation}
物理通量取为
\begin{equation}
  \F(\U) =
  \begin{pmatrix}
    hu \\[2pt]
    \dfrac{(hu)^2}{h} + \dfrac{1}{2}gh^2 \\[8pt]
    \dfrac{hu\,hv}{h}
  \end{pmatrix},
  \qquad
  \G(\U) =
  \begin{pmatrix}
    hv \\[2pt]
    \dfrac{hu\,hv}{h} \\[8pt]
    \dfrac{(hv)^2}{h} + \dfrac{1}{2}gh^2
  \end{pmatrix},
  \label{eq:fluxes}
\end{equation}
其中 $g = 9.81\,\mathrm{m\,s^{-2}}$ 为重力加速度。

\subsection{特征结构}

方程~\eqref{eq:swe} 构成双曲型守恒律组。$x$ 方向通量 $\F$ 的 Jacobian 矩阵为
\begin{equation}
  A(\U) = \frac{\partial\F}{\partial\U}
  =
  \begin{pmatrix}
    0 & 1 & 0 \\
    -u^2 + c^2 & 2u & 0 \\
    -uv & v & u
  \end{pmatrix},
  \label{eq:jacobian-x}
\end{equation}
其中 $c = \sqrt{gh}$ 为重力波波速。矩阵 $A$ 的三个特征值为
\begin{equation}
  \lambda_1 = u - c,\qquad
  \lambda_2 = u,\qquad
  \lambda_3 = u + c,
  \label{eq:eigen-x}
\end{equation}
对应的右特征向量为
\begin{equation}
  \bm{r}_1 = \begin{pmatrix}1 & u-c & v\end{pmatrix}\T,\quad
  \bm{r}_2 = \begin{pmatrix}0 & 0 & 1\end{pmatrix}\T,\quad
  \bm{r}_3 = \begin{pmatrix}1 & u+c & v\end{pmatrix}\T.
  \label{eq:eigenvec-x}
\end{equation}
特征值 $u\pm c$ 对应沿 $x$ 方向传播的重力波，$\lambda_2 = u$ 对应随流输运的
涡旋波模态。$y$ 方向通量 $\G$ 的 Jacobian 矩阵 $B(\U) = \partial\G/\partial\U$
具有类似的结构，其特征值为
\begin{equation}
  \mu_1 = v - c,\qquad
  \mu_2 = v,\qquad
  \mu_3 = v + c.
  \label{eq:eigen-y}
\end{equation}
谱半径 $\rho(A) = |u| + c$ 决定了 Rusanov 通量中数值耗散的强度，也是 CFL
时间步长约束的理论基础。

\subsection{有限体积离散}

计算域取为矩形区域 $[0, L_x] \times [0, L_y]$，以 $N_x \times N_y$ 个均匀笛卡尔
网格单元进行离散。单元中心坐标为
\[
  x_i = \Bigl(i + \tfrac{1}{2}\Bigr)\Delta x,\qquad
  y_j = \Bigl(j + \tfrac{1}{2}\Bigr)\Delta y,
\]
其中 $\Delta x = L_x/N_x$，$\Delta y = L_y/N_y$。将方程~\eqref{eq:swe} 在
每个控制体上积分，得到半离散有限体积格式
\begin{equation}
  \frac{d\U_{i,j}}{dt}
  =
  -\frac{\widehat{\F}_{i+1/2,j} - \widehat{\F}_{i-1/2,j}}{\Delta x}
  -\frac{\widehat{\G}_{i,j+1/2} - \widehat{\G}_{i,j-1/2}}{\Delta y}
  \equiv \mathcal{L}(\U)_{i,j},
  \label{eq:semidiscrete}
\end{equation}
其中 $\widehat{\F}$ 和 $\widehat{\G}$ 分别为 $x$ 和 $y$ 方向的数值通量。

边界条件取零梯度外推：当通量计算需要域外虚拟单元时，直接复制最近内部单元的
状态。这一处理简单高效，适用于短时间合成实验。

\subsection{Rusanov 数值通量}

界面数值通量采用 Rusanov（局部 Lax--Friedrichs）格式。以 $x$ 方向界面为例，
设界面左、右两侧状态分别为 $\U_L$ 和 $\U_R$，Rusanov 通量定义为
\begin{equation}
  \widehat{\F}(\U_L, \U_R)
  =
  \frac{1}{2}\bigl[\F(\U_L) + \F(\U_R)\bigr]
  - \frac{1}{2}s_x(\U_R - \U_L),
  \label{eq:rusanov-x}
\end{equation}
其中局部最大信号速度为
\begin{equation}
  s_x = \max\bigl(|u_L| + \sqrt{gh_L},\; |u_R| + \sqrt{gh_R}\bigr).
  \label{eq:smax}
\end{equation}
$y$ 方向通量具有完全类似的形式。Rusanov 通量的第一项为中心通量，第二项为标量
数值耗散项。它不依赖完整的特征分解，实现简单且数值鲁棒性良好。

\subsection{时间积分与 CFL 条件}

记~\eqref{eq:semidiscrete} 中的空间离散算子为 $\mathcal{L}$。时间推进采用
Heun 型二阶 Runge--Kutta 格式：
\begin{align}
  \U^{(1)} &= \U^n + \Delta t\,\mathcal{L}(\U^n), \label{eq:rk2-stage1} \\
  \U^{(2)} &= \U^{(1)} + \Delta t\,\mathcal{L}(\U^{(1)}), \label{eq:rk2-stage2} \\
  \U^{n+1} &= \frac{1}{2}\U^n + \frac{1}{2}\U^{(2)}. \label{eq:rk2-final}
\end{align}

时间步长 $\Delta t$ 由 CFL 条件自适应确定：
\begin{equation}
  \Delta t
  = \mathrm{CFL} \cdot
  \min_{i,j}\left(
    \frac{\Delta x}{|u_{i,j}| + \sqrt{gh_{i,j}}},\;
    \frac{\Delta y}{|v_{i,j}| + \sqrt{gh_{i,j}}}
  \right),
  \label{eq:cfl}
\end{equation}
其中 CFL 数通常取为 $0.40$。

为保障数值稳定性，所有计算中对水深施加最小阈值 $h_{\min} > 0$。当某个单元的
水深因更新而低于该阈值时，将其重置为 $h = h_{\min}$，并令该点的动量分量为零。
这一约束不属于浅水方程本身，而是为消除近干单元的数值奇性而引入的物理可接受性
保护。

%=====================================================================
\section{集合数据同化框架}
\label{sec:framework}
%=====================================================================

\subsection{状态向量与观测算子}

从数据同化的角度看，单一集合成员的状态由所有网格点上的守恒变量拼接而成：
\begin{equation}
  \x
  =
  \begin{pmatrix}
    h_1, \dots, h_N,\;
    (hu)_1, \dots, (hu)_N,\;
    (hv)_1, \dots, (hv)_N
  \end{pmatrix}\T
  \in \R^{3N},
  \label{eq:full-state}
\end{equation}
其中 $N = N_x N_y$ 为网格单元总数。因此，尽管前向模式自然地将三个物理分量
分别存储，从滤波器的角度看，同化的对象是包含 $h$、$hu$ 和 $hv$ 的完整状态
向量。

设同化系统的观测为在 $m$ 个选定网格点上抽取的水深值。记第 $p$ 个观测对应的
网格点索引为 $k_p$，观测算子 $\calH: \R^{3N} \to \R^m$ 定义为
\begin{equation}
  [\calH(\x)]_p = h_{k_p}, \qquad p = 1, 2, \dots, m.
  \label{eq:obs-operator}
\end{equation}
观测方程为
\begin{equation}
  \y = \calH(\x^{\mathrm{truth}}) + \bm{\varepsilon},
  \qquad
  \bm{\varepsilon} \sim \mathcal{N}(0, R),
  \label{eq:obs-model}
\end{equation}
其中 $\x^{\mathrm{truth}}$ 为未知的真实状态。在本文的合成实验中，观测误差
协方差取为对角形式
\begin{equation}
  R = \sigma_{\mathrm{obs}}^2 I_m,
  \label{eq:R-diag}
\end{equation}
即不同观测位置的误差相互独立且具有相同的方差 $\sigma_{\mathrm{obs}}^2$。

需要指出的是，虽然观测算子仅抽取水深分量，但 Kalman 增益中包含了完整状态
与观测之间的互协方差。因此，分析更新不仅修正被直接观测的 $h$，也会通过
样本相关性同步修正未被观测的动量分量 $hu$ 和 $hv$。这正是集合方法的
核心优势之一。

\subsection{集合初始化}

真值场的初始条件取为静止水面上叠加一个高斯型水面隆起：
\begin{equation}
  h(x, y, 0) = 1 + 0.2\,
  \exp\!\left(-\frac{(x - x_0)^2 + (y - y_0)^2}{0.02}\right),
  \qquad hu = 0,\quad hv = 0.
  \label{eq:truth-init}
\end{equation}

初始集合成员由真值初值添加高斯随机扰动生成：
\begin{align}
  h^{(e)} &= h^{\mathrm{truth}} + \sigma_{h0}\,\xi_h^{(e)}, \label{eq:init-h} \\
  (hu)^{(e)} &= (hu)^{\mathrm{truth}} + \sigma_{m0}\,\xi_{hu}^{(e)}, \label{eq:init-hu} \\
  (hv)^{(e)} &= (hv)^{\mathrm{truth}} + \sigma_{m0}\,\xi_{hv}^{(e)}, \label{eq:init-hv}
\end{align}
其中 $\xi_h^{(e)}, \xi_{hu}^{(e)}, \xi_{hv}^{(e)}$ 为独立标准正态随机变量，
$e = 1, \dots, N_e$。典型参数取值 $\sigma_{h0} = 0.02$，$\sigma_{m0} = 0.005$。
记初始分析集合为 $\{\x_0^{a,(e)}\}_{e=1}^{N_e}$。

\subsection{预报步与样本统计量}

在第 $k$ 个同化窗口，设前一个分析时刻的集合为 $\{\x_{k-1}^{a,(e)}\}$。
预报步将每个成员通过非线性浅水模式向前积分一个同化窗长 $\tau$：
\begin{equation}
  \x_k^{f,(e)} = \calM_\tau\bigl(\x_{k-1}^{a,(e)}\bigr),
  \qquad e = 1, 2, \dots, N_e,
  \label{eq:forecast}
\end{equation}
其中 $\calM_\tau$ 表示由第~\ref{sec:forward} 节所述数值模式定义的
非线性预报算子，$\tau$ 为同化窗的物理时长。

预报集合的样本均值与归一化扰动矩阵定义为
\begin{equation}
  \bar{\x}^f = \frac{1}{N_e}\sum_{e=1}^{N_e} \x^{f,(e)},
  \qquad
  X^f = \frac{1}{\sqrt{N_e - 1}}
  \begin{bmatrix}
    \x^{f,(1)} - \bar{\x}^f & \cdots & \x^{f,(N_e)} - \bar{\x}^f
  \end{bmatrix}.
  \label{eq:ensemble-stat}
\end{equation}
于是，预报误差协方差的样本近似为
\begin{equation}
  P^f \approx X^f (X^f)\T.
  \label{eq:Pf-approx}
\end{equation}

将每个预报成员映射到观测空间：
\begin{equation}
  \y^{f,(e)} = \calH(\x^{f,(e)}), \qquad e = 1, \dots, N_e.
  \label{eq:obs-mapping}
\end{equation}
观测空间的预报均值为
\[
  \bar{\y}^f = \frac{1}{N_e}\sum_{e=1}^{N_e} \y^{f,(e)}.
\]
对应的归一化观测扰动矩阵为
\begin{equation}
  Y^f = \frac{1}{\sqrt{N_e - 1}}
  \begin{bmatrix}
    \y^{f,(1)} - \bar{\y}^f & \cdots & \y^{f,(N_e)} - \bar{\y}^f
  \end{bmatrix}.
  \label{eq:Yf}
\end{equation}
对于线性观测算子 $H$（即本文的水深抽取算子），有 $Y^f = H X^f$。对一般的
非线性观测算子，$Y^f$ 以逐成员映射后减均值的方式近似，等价于以集合差分
替代了切线性观测算子的作用 \cite{asch2016data}。

%=====================================================================
\section{随机集合卡尔曼滤波}
\label{sec:enkf}
%=====================================================================

\subsection{从 Kalman 滤波到集合近似}

考虑线性高斯观测系统 $\y = H\x + \bm{\varepsilon}$，其中
$\bm{\varepsilon} \sim \mathcal{N}(0, R)$。标准 Kalman 滤波的分析均值和
协方差更新公式为
\begin{align}
  \x^a &= \x^f + K(\y - H\x^f), \label{eq:kf-mean} \\
  P^a &= (I - KH)P^f, \label{eq:kf-cov}
\end{align}
其中 Kalman 增益为
\begin{equation}
  K = P^f H\T (H P^f H\T + R)^{-1}.
  \label{eq:kf-gain}
\end{equation}

集合卡尔曼滤波的核心思想是用有限集合成员的样本统计量代替真实的 $P^f$ 和
$H P^f H\T$。定义归一化预报扰动矩阵 $X^f$ 和观测扰动矩阵 $Y^f$ 如
\eqref{eq:ensemble-stat}--\eqref{eq:Yf}，则样本协方差近似为
\begin{equation}
  P_{xy} = X^f (Y^f)\T \approx P^f H\T,\qquad
  P_{yy} = Y^f (Y^f)\T \approx H P^f H\T.
  \label{eq:sample-cov}
\end{equation}

若直接将确定性分析公式 $\x^{a,(e)} = \x^{f,(e)} + K(\y - H\x^{f,(e)})$
施加于每个集合成员，则分析离差满足 $X^a = (I - KH)X^f$，样本分析协方差
变为 $P^a_{\mathrm{sample}} = (I - KH)P^f(I - KH)\T$。与标准 Kalman
分析协方差~\eqref{eq:kf-cov} 相比，该结果缺失了 $K R K\T$ 项，从而导致
集合离散度系统性偏低。

\subsection{扰动观测分析更新}

为补足缺失的 $K R K\T$ 项，随机 EnKF \cite{burgers1998analysis} 对每个
集合成员生成独立的扰动观测：
\begin{equation}
  \y^{(e)}_{\mathrm{pert}} = \y + \bm{u}^{(e)},
  \qquad \bm{u}^{(e)} \sim \mathcal{N}(0, R),\quad
  e = 1, \dots, N_e.
  \label{eq:pert-obs}
\end{equation}
随后对每个成员执行分析更新：
\begin{equation}
  \x^{a,(e)} = \x^{f,(e)} + K\bigl(\y^{(e)}_{\mathrm{pert}} - \calH(\x^{f,(e)})\bigr).
  \label{eq:enkf-update}
\end{equation}

在这一构造下，分析离差除 $(I - KH)X^f$ 外，还包含观测扰动 $\bm{u}^{(e)}$
诱导的随机分量。对该随机分量取期望后，缺失的 $K R K\T$ 项得以恢复，
从而使得
\[
  \mathbb{E}\bigl[P^a_{\mathrm{sample}}\bigr]
  = (I - KH)P^f(I - KH)\T + K R K\T
  = (I - KH)P^f.
\]
因此，扰动观测并非任意的数值技巧，而是保证有限集合分析协方差在期望意义上
与 Kalman 理论一致的数学构造。

\subsection{样本协方差与 Kalman 增益的计算}

将 $P^f \approx X^f(X^f)\T$ 代入~\eqref{eq:kf-gain}，可将 Kalman 增益改写为
以样本协方差表示的形式：
\begin{equation}
  K = P_{xy}(P_{yy} + R)^{-1},
  \label{eq:enkf-K}
\end{equation}
其中 $P_{xy}$ 和 $P_{yy}$ 由~\eqref{eq:sample-cov} 定义。当观测误差协方差取
对角形式 $R = \sigma_{\mathrm{obs}}^2 I$ 时，$P_{yy} + R$ 仅为 $m \times m$
矩阵，其求逆代价随观测维数 $m$ 线性增长。

样本协方差的构造采用标准无偏估计：
\begin{align}
  P_{xy} &= \frac{1}{N_e - 1}
    \sum_{e=1}^{N_e}
    \bigl(\x^{f,(e)} - \bar{\x}^f\bigr)
    \bigl(\y^{f,(e)} - \bar{\y}^f\bigr)\T, \label{eq:Pxy-sum} \\
  P_{yy} &= \frac{1}{N_e - 1}
    \sum_{e=1}^{N_e}
    \bigl(\y^{f,(e)} - \bar{\y}^f\bigr)
    \bigl(\y^{f,(e)} - \bar{\y}^f\bigr)\T. \label{eq:Pyy-sum}
\end{align}

\subsection{算法小结}

随机 EnKF 在一个同化窗口内的完整分析流程可概括如下：

\begin{enumerate}[label=(\arabic*), leftmargin=*]
  \item 由预报集合计算样本均值 $\bar{\x}^f$ 和 $\bar{\y}^f$；
  \item 按~\eqref{eq:Pxy-sum}--\eqref{eq:Pyy-sum} 计算样本协方差 $P_{xy}$ 和 $P_{yy}$；
  \item 构造 $S = P_{yy} + R$，求逆得到 $S^{-1}$，计算 Kalman 增益 $K = P_{xy}S^{-1}$；
  \item 对每个成员生成扰动观测 $\y^{(e)}_{\mathrm{pert}} = \y + \bm{u}^{(e)}$，其中
        $\bm{u}^{(e)} \sim \mathcal{N}(0, R)$；
  \item 按~\eqref{eq:enkf-update} 更新每个成员，得到分析集合 $\{\x^{a,(e)}\}$。
\end{enumerate}

分析完成后，对每个成员施加物理约束：若某单元水深低于最小阈值 $h_{\min}$，
则将其重置为 $h = h_{\min}$ 并令动量分量为零。该约束属于数值鲁棒性保护，
并非 EnKF 理论自身的组成部分。

%=====================================================================
\section{确定性集合平方根滤波}
\label{sec:ensrf}
%=====================================================================

\subsection{均值--扰动分离}

随机 EnKF 通过在观测上添加扰动来实现协方差的正确收缩，但这一构造同时引入
了额外的采样噪声。集合平方根滤波（EnSRF）是确定性集合滤波的一类代表方法，
其核心思想是：\textbf{不扰动观测，而是将分析步分解为均值更新与扰动变换两个
相互独立的部分，并对扰动施加确定性的平方根收缩}。

给定预报集合的均值 $\bar{\x}^f$ 和归一化扰动矩阵 $X^f$，均值更新仍沿用
Kalman 公式：
\begin{equation}
  \bar{\x}^a = \bar{\x}^f + K\bigl(\y - \calH(\bar{\x}^f)\bigr),
  \label{eq:ensrf-mean}
\end{equation}
其中 $K$ 由~\eqref{eq:enkf-K} 给出。

扰动更新的目标是寻找一个 $N_e \times N_e$ 的变换矩阵 $T$，使得
\begin{equation}
  X^a = X^f T,
  \label{eq:ensrf-pert}
\end{equation}
且满足 Kalman 分析协方差关系：
\begin{equation}
  X^a (X^a)\T = P^a = (I - KH)P^f.
  \label{eq:ensrf-consistency}
\end{equation}
将 $P^f \approx X^f (X^f)\T$ 代入，可得 $T$ 必须满足的约束条件：
\begin{equation}
  T T\T = I - (Y^f)\T (P_{yy} + R)^{-1} Y^f.
  \label{eq:T-constraint}
\end{equation}

满足~\eqref{eq:T-constraint} 的 $T$ 存在无穷多种选取方式。不同的选取对应
文献中不同的确定性集合滤波变体。因此，所有平方根滤波方法在数学上都可以统一
为\textbf{求一个满足 $T T\T = I - (Y^f)\T(P_{yy}+R)^{-1}Y^f$ 的变换矩阵
$T$，并以 $X^a = X^f T$ 更新集合扰动}这一框架。

\subsection{集合子空间公式}

利用 Sherman--Morrison--Woodbury 恒等式，可将 $T$ 的约束条件改写为更紧凑
的集合子空间形式。定义 $N_e \times N_e$ 对称正定矩阵
\begin{equation}
  T_e = I + (Y^f)\T R^{-1} Y^f,
  \label{eq:Te}
\end{equation}
则约束~\eqref{eq:T-constraint} 等价于
\begin{equation}
  T T\T = T_e^{-1}.
  \label{eq:T-rel}
\end{equation}

均值更新同样可通过集合子空间表达。将 $K$ 的表达式代入~\eqref{eq:ensrf-mean}，
经代数整理后可得
\begin{equation}
  \bar{\x}^a = \bar{\x}^f + X^f \w^a,
  \label{eq:mean-weight}
\end{equation}
其中权重向量 $\w^a \in \R^{N_e}$ 满足低维线性系统
\begin{equation}
  \w^a = T_e^{-1} (Y^f)\T R^{-1} \bigl(\y - \calH(\bar{\x}^f)\bigr).
  \label{eq:wa}
\end{equation}

在~\eqref{eq:T-rel}--\eqref{eq:wa} 的表述下，确定性集合滤波的计算归结为
两个 $N_e \times N_e$ 矩阵运算：求解 $T_e \w^a = (Y^f)\T R^{-1}\mathbf{d}$
以获得均值权重，以及从 $T_e^{-1}$ 中提取一个平方根因子 $T$ 以更新扰动。
当 $N_e \ll \dim(\x)$ 时（这是集合方法的典型工况），上述运算的计算代价
远低于在完整状态空间中操作协方差矩阵。

\subsection{对称平方根变换}

最直接的平方根选取是令 $T$ 为 $T_e^{-1}$ 的对称平方根：
\begin{equation}
  T = T_e^{-1/2}.
  \label{eq:sym-sqrt}
\end{equation}
实现时，对 $T_e$ 进行特征分解：
\begin{equation}
  T_e = V \Lambda V\T,
  \label{eq:eigendecomp}
\end{equation}
其中 $V$ 为特征向量矩阵，$\Lambda = \operatorname{diag}(\lambda_1, \dots,
\lambda_{N_e})$ 为特征值对角阵（全部为正，因 $T_e$ 对称正定）。则
\begin{equation}
  T = V \Lambda^{-1/2} V\T,
  \qquad
  \Lambda^{-1/2} = \operatorname{diag}(\lambda_1^{-1/2}, \dots,
  \lambda_{N_e}^{-1/2}).
  \label{eq:T-symsqrt}
\end{equation}

对称平方根具有以下性质：(i)~$T$ 本身为对称矩阵，变换后集合均值不受扰动
变换的偏置影响；(ii)~$T$ 由 $T_e$ 唯一确定，不需要额外的参数选择；
(iii)~$T$ 使分析集合扰动保持无偏。

\subsection{分析集合重构}

综合均值更新~\eqref{eq:mean-weight} 和扰动变换~\eqref{eq:ensrf-pert}，
第 $e$ 个分析成员的重构公式为
\begin{equation}
  \x^{a,(e)}
  = \bar{\x}^f
  + X^f \w^a
  + \sqrt{N_e - 1}\; X^f [T]_e,
  \label{eq:ensrf-reconstruct}
\end{equation}
其中 $[T]_e$ 表示矩阵 $T$ 的第 $e$ 列。该式同时更新 $h$、$hu$ 和 $hv$ 三个
物理分量，观测信息通过 $Y^f$ 中的异常相关性传播至所有状态分量。

算法~\ref{alg:ensrf} 给出了基于对称平方根的 EnSRF 在一个同化窗口内的完整流程。

\begin{algorithm}[htbp]
  \caption{对称平方根 EnSRF 分析步}
  \label{alg:ensrf}
  \begin{algorithmic}[1]
    \Require 预报集合 $\{\x^{f,(e)}\}_{e=1}^{N_e}$，观测 $\y$，
             观测误差协方差 $R$
    \Ensure 分析集合 $\{\x^{a,(e)}\}_{e=1}^{N_e}$

    \State 计算预报均值 $\bar{\x}^f = \frac{1}{N_e}\sum_e \x^{f,(e)}$，
           $\bar{\y}^f = \frac{1}{N_e}\sum_e \calH(\x^{f,(e)})$

    \State 构造归一化扰动矩阵 $X^f$ 和 $Y^f$（式~\eqref{eq:ensemble-stat}、
           \eqref{eq:Yf}）

    \State 构造创新向量 $\mathbf{d} = \y - \bar{\y}^f$

    \State 构造 $N_e \times N_e$ 矩阵 $T_e = I + (Y^f)\T R^{-1} Y^f$

    \State 求均值权重：$\w^a = T_e^{-1} (Y^f)\T R^{-1} \mathbf{d}$

    \State 对 $T_e$ 做特征分解 $T_e = V \Lambda V\T$

    \State 构造对称平方根 $T = V \Lambda^{-1/2} V\T$

    \State 重构分析成员 $\x^{a,(e)} = \bar{\x}^f + X^f\w^a
           + \sqrt{N_e-1}\,X^f[T]_e$，$e = 1,\dots,N_e$

    \State 施加重力约束：若 $h_{i,j} < h_{\min}$，则
           $h_{i,j} \leftarrow h_{\min}$，$(hu)_{i,j} \leftarrow 0$，
           $(hv)_{i,j} \leftarrow 0$
  \end{algorithmic}
\end{algorithm}

\subsection{与随机 EnKF 的比较}

从理论上说，随机 EnKF 与确定性 EnSRF 的唯一区别在于如何使分析集合协方差
满足 $P^a = (I - KH)P^f$：
\begin{itemize}
  \item 随机 EnKF 在观测上添加随机扰动，使分析离差的期望协方差包含 $K R K\T$
        项，从而在期望意义上满足该关系；
  \item 确定性 EnSRF 直接构造 $X^a$ 使得 $X^a(X^a)\T$ 精确满足该关系，
        不含观测扰动引入的额外采样噪声。
\end{itemize}

在有限集合规模下，确定性方案的分析误差和离散度曲线通常更平滑，但其代价
是需要额外的 $N_e \times N_e$ 矩阵特征分解或 Cholesky 分解。当 $N_e$ 较小时，
这一开销可以忽略。两类方法在非线性较强或集合离散度不足时，都可能需要引入
协方差膨胀（covariance inflation）和局地化（localization）以维持合理的
集合散布。

%=====================================================================
\section{数值实验}
\label{sec:experiments}
%=====================================================================

\subsection{实验配置}

为定量比较随机 EnKF 和确定性 EnSRF 的同化性能，设计如下双系统合成实验。
计算域取为 $[0, 1] \times [0, 1]$，网格分辨率 $N_x = N_y = 50$。集合规模
$N_e = 20$，同化窗长 $\tau = 0.04$ 个无量纲时间单位，共进行 31 个同化循环。
观测网络在计算域内以 $5 \times 5$ 的均匀分布选取 25 个水深观测点，观测误差
标准差 $\sigma_{\mathrm{obs}} = 0.01$。两组滤波器共享相同的真值轨道、观测
数据、初始集合和模式参数，仅分析步算法不同：随机 EnKF 使用扰动观测方案
\eqref{eq:enkf-update}，EnSRF 使用对称平方根方案（算法~\ref{alg:ensrf}）。

\subsection{诊断量}

定义第 $k$ 个同化窗口内水深分量的均方根误差为
\begin{equation}
  \mathrm{RMSE}_k = \sqrt{
    \frac{1}{N}\sum_{i,j}
    \bigl(h_{i,j}^{\mathrm{mean}} - h_{i,j}^{\mathrm{truth}}\bigr)^2
  },
  \label{eq:rmse}
\end{equation}
其中 $h^{\mathrm{mean}}$ 为分析均值（或预测均值）的水深场。分别对两种滤波器
计算预测 RMSE（$\mathrm{RMSE}_f$）和分析 RMSE（$\mathrm{RMSE}_a$），以检验
观测信息是否有效降低了预报误差。进一步，定义集合离散度
\[
  \mathrm{spread}_k = \sqrt{
    \frac{1}{N}\sum_{i,j} \operatorname{var}\bigl(\{h_{i,j}^{(e)}\}_{e=1}^{N_e}\bigr)
  },
\]
以及预测创新均方根
\[
  \mathrm{innov}_k = \sqrt{
    \frac{1}{m}\sum_{p=1}^{m}
    \bigl(y_p - \bar{y}_p^f\bigr)^2
  },
\]
用以监测滤波器的一致性。

\subsection{结果与讨论}

\begin{figure}[htbp]
  \centering
  \IfFileExists{comparison_output/validation_metrics.png}{%
    \includegraphics[width=\textwidth]{comparison_output/validation_metrics.png}%
  }{%
    \fbox{\parbox{0.95\textwidth}{\centering\vspace{2em}
      \textbf{[图~\ref{fig:validation-metrics}]}\\[1ex]
      \texttt{comparison\_output/validation\_metrics.png}\\[0.5ex]
      \footnotesize 运行 \texttt{./da \&\& python3 plot\_validation_figures.py} 生成此图
      \vspace{2em}}}%
  }
  \caption{双系统同化实验的整体统计量。(a) 预测/分析 RMSE 时间序列；
           (b) 预测/分析 spread；(c) 预测创新 RMS；(d) EnSRF 与随机 EnKF 的
           分析 RMSE 比值，低于 1 表示 EnSRF 在该循环中更优。}
  \label{fig:validation-metrics}
\end{figure}

图~\ref{fig:validation-metrics} 给出了四个诊断量的完整时间序列。主要发现如下：

\begin{enumerate}[leftmargin=*]
  \item 在 31 个同化循环中，随机 EnKF 有 15 个循环满足
        $\mathrm{RMSE}_a < \mathrm{RMSE}_f$，EnSRF 有 16 个循环满足该条件。
        两种方法均能在大约半数的循环中将观测信息转化为净误差下降；
  \item EnSRF 在 31 个循环中的 21 个循环里取得了更小的分析 RMSE。全时段
        平均而言，$\mathrm{RMSE}_{a,\mathrm{EnSRF}} /
        \mathrm{RMSE}_{a,\mathrm{EnKF}} \approx 0.934$，即 EnSRF 的平均
        分析误差较随机 EnKF 低约 6.6\%；
  \item 两种方法的创新 RMS 曲线高度吻合，表明它们吸收的是同一批观测信息，
        性能差异主要来源于分析步内部对集合离差的不同处理方式；
  \item 在整个实验期间未触发物理裁剪（即分析更新后无水深刻度低于
        $h_{\min}$），说明两种滤波器在同化过程中均保持了物理可接受性。
\end{enumerate}

\begin{figure}[htbp]
  \centering
  \IfFileExists{comparison_output/validation_snapshots.png}{%
    \includegraphics[width=\textwidth]{comparison_output/validation_snapshots.png}%
  }{%
    \fbox{\parbox{0.95\textwidth}{\centering\vspace{2em}
      \textbf{[图~\ref{fig:validation-snapshots}]}\\[1ex]
      \texttt{comparison\_output/validation\_snapshots.png}\\[0.5ex]
      \footnotesize 运行 \texttt{./da \&\& python3 plot\_validation_figures.py} 生成此图
      \vspace{2em}}}%
  }
  \caption{代表性循环的场误差快照。每列对应一个同化循环（初期、中期、末期）；
           第一行为真值场，第二行为随机 EnKF 分析误差，第三行为 EnSRF 分析误差；
           加号标记观测位置。}
  \label{fig:validation-snapshots}
\end{figure}

图~\ref{fig:validation-snapshots} 选取了初期、中期和末期三个代表性循环的
误差空间分布。在初期循环，两种方法均迅速将误差压缩至观测网络附近的局部尺度。
到中后期，大范围系统性偏差已基本消除，剩余误差主要表现为局部弱振荡。在末期
循环，EnSRF 的误差幅值普遍略小于随机 EnKF，与 RMSE 时间序列的统计结论一致。

综合上述结果，在 $N_e = 20$ 的有限集合规模下，确定性 EnSRF 在分析精度上
整体略优于随机 EnKF，验证了平方根变换消除额外采样噪声的理论预期。两类方法
均为 $h$ 观测在浅水方程状态空间中提供了有效的同化手段。

%=====================================================================
\section{总结与讨论}
\label{sec:conclusion}
%=====================================================================

本文从 Kalman 滤波的数学基础出发，系统阐述了应用于二维浅水方程的随机 EnKF
和确定性 EnSRF 方法的理论基础。前向模式采用守恒型浅水方程的有限体积离散，
空间通量为 Rusanov 格式，时间推进为二阶 Runge--Kutta 格式。在同化框架中，
状态向量包含水深和两个动量分量，观测算子从网格点抽取水深值。

随机 EnKF 通过样本协方差近似 Kalman 增益，并在观测上添加随机扰动以保证分析
协方差的统计一致性。确定性 EnSRF 将分析步分解为均值更新和扰动变换，利用对称
平方根对集合离差施加确定性的收缩变换。两类方法共享同一个集合子空间框架，其
差异仅在于如何实现 $P^a = (I - KH)P^f$ 这一关系。串行 EnSRF、ETKF 和 EAKF
等变体均可纳入 $T T\T = T_e^{-1}$ 的统一数学结构。

双系统合成实验的结果表明，在有限集合规模下，确定性 EnSRF 的整体表现略优于
随机 EnKF，其分析误差平均降低约 6.6\%，且误差曲线更为平滑。值得指出的是，
本文的数值实验未包含协方差膨胀和局地化等附加技术。在高维或强非线性问题中，
这些技术通常是维持集合合理离散度和抑制远距离虚假相关的必要手段，也是进一步
研究方向。

\appendix
\section{EnKF 分析步的详细算法描述}
\label{app:enkf}

本附录给出随机 EnKF 在一个同化窗口内的逐步数学描述，与主文
第~\ref{sec:enkf} 节互为补充。

\medskip\noindent
\textbf{输入：}预报集合 $\{\x^{f,(e)}\}_{e=1}^{N_e}$，其中每个成员的状态
由 $h$、$hu$、$hv$ 三部分组成；观测向量 $\y \in \R^m$；观测索引
$\{k_p\}_{p=1}^{m}$，其中 $k_p$ 为第 $p$ 个观测对应的网格点编号；观测误差
标准差 $\sigma_{\mathrm{obs}}$。

\medskip\noindent
\textbf{输出：}分析集合 $\{\x^{a,(e)}\}_{e=1}^{N_e}$。

\medskip\noindent
\textbf{步骤 1：计算预报集合均值。}
对每个状态分量在所有成员上求算术平均：
\[
  \bar{\x}^f = \frac{1}{N_e}\sum_{e=1}^{N_e} \x^{f,(e)}.
\]

\medskip\noindent
\textbf{步骤 2：映射至观测空间。}
对每个成员 $e$ 和每个观测点 $p$：
\[
  y_p^{f,(e)} = \calH(\x^{f,(e)})_p = h_{k_p}^{f,(e)}.
\]
计算观测空间均值：
\[
  \bar{\y}^f = \frac{1}{N_e}\sum_{e=1}^{N_e} \y^{f,(e)}.
\]

\medskip\noindent
\textbf{步骤 3：构造样本协方差矩阵。}
对每个成员 $e$ 和每个观测分量 $p$，计算观测离差
\[
  \delta y_p^{(e)} = y_p^{f,(e)} - \bar{y}_p^f.
\]
对每个状态分量 $i$（依次遍历 $h$、$hu$、$hv$），计算状态离差
\[
  \delta x_i^{(e)} = x_i^{f,(e)} - \bar{x}_i^f,
\]
并按无偏估计累加互协方差和自协方差：
\begin{align*}
  (P_{xy})_{i,p} &= \frac{1}{N_e-1}\sum_{e=1}^{N_e} \delta x_i^{(e)}\,\delta y_p^{(e)}, \\[4pt]
  (P_{yy})_{p,q} &= \frac{1}{N_e-1}\sum_{e=1}^{N_e} \delta y_p^{(e)}\,\delta y_q^{(e)}.
\end{align*}

\medskip\noindent
\textbf{步骤 4：计算 Kalman 增益。}
构造 $S = P_{yy} + \sigma_{\mathrm{obs}}^2 I_m$，求逆得 $S^{-1}$，计算
\[
  K = P_{xy} S^{-1} \;\in\; \R^{3N \times m}.
\]
矩阵 $K$ 的第 $i$ 行第 $p$ 列量化了第 $p$ 个观测对第 $i$ 个状态分量的修正权重。

\medskip\noindent
\textbf{步骤 5：生成扰动观测与创新。}
对每个成员 $e$ 和每个观测分量 $p$，生成扰动观测
\[
  y_{\mathrm{pert},p}^{(e)} = y_p + \sigma_{\mathrm{obs}}\,\eta_p^{(e)},
  \qquad \eta_p^{(e)} \sim \mathcal{N}(0,1),
\]
并计算创新向量
\[
  \mathbf{d}^{(e)} = \y_{\mathrm{pert}}^{(e)} - \y^{f,(e)}.
\]

\medskip\noindent
\textbf{步骤 6：分析更新。}
对每个成员 $e$ 和每个状态分量 $i$，计算增量并更新：
\[
  x_i^{a,(e)} = x_i^{f,(e)} + \sum_{p=1}^{m} K_{i,p}\,d_p^{(e)}.
\]
以向量形式即
\[
  \x^{a,(e)} = \x^{f,(e)} + K\,\mathbf{d}^{(e)}.
\]

\medskip\noindent
\textbf{步骤 7：物理约束。}
对每个成员 $e$ 和每个网格单元 $(i,j)$，若 $h_{i,j}^{a,(e)} < h_{\min}$，则
\[
  h_{i,j}^{a,(e)} \leftarrow h_{\min},\qquad
  (hu)_{i,j}^{a,(e)} \leftarrow 0,\qquad
  (hv)_{i,j}^{a,(e)} \leftarrow 0.
\]

\section{EnSRF 分析步的详细算法描述}
\label{app:ensrf}

本附录给出对称平方根 EnSRF 在一个同化窗口内的逐步数学描述，与主文
第~\ref{sec:ensrf} 节互为补充。算法流程与算法~\ref{alg:ensrf} 一致，
此处展开为逐矩阵操作的详细说明。

\medskip\noindent
\textbf{输入：}预报集合 $\{\x^{f,(e)}\}_{e=1}^{N_e}$；观测向量 $\y \in \R^m$；
观测索引 $\{k_p\}_{p=1}^{m}$；观测误差标准差 $\sigma_{\mathrm{obs}}$。

\medskip\noindent
\textbf{输出：}分析集合 $\{\x^{a,(e)}\}_{e=1}^{N_e}$。

\medskip\noindent
\textbf{步骤 1：计算预报均值。}
\begin{align*}
  \bar{\x}^f &= \frac{1}{N_e}\sum_{e=1}^{N_e} \x^{f,(e)}, \\
  \bar{\y}^f &= \frac{1}{N_e}\sum_{e=1}^{N_e} \calH(\x^{f,(e)}).
\end{align*}

\medskip\noindent
\textbf{步骤 2：构造异常矩阵。}
保存原始状态异常矩阵（不除以 $\sqrt{N_e-1}$）：
\[
  A = \begin{bmatrix}
    \x^{f,(1)} - \bar{\x}^f & \cdots & \x^{f,(N_e)} - \bar{\x}^f
  \end{bmatrix} \in \R^{3N \times N_e}.
\]
构造归一化观测异常矩阵：
\[
  Y^f = \frac{1}{\sqrt{N_e-1}}
  \begin{bmatrix}
    \y^{f,(1)} - \bar{\y}^f & \cdots & \y^{f,(N_e)} - \bar{\y}^f
  \end{bmatrix} \in \R^{m \times N_e}.
\]
注意 $A = \sqrt{N_e-1}\,X^f$，其中 $X^f$ 为归一化扰动矩阵。

\medskip\noindent
\textbf{步骤 3：构造集合子空间矩阵。}
\[
  T_e = I_{N_e} + (Y^f)\T R^{-1} Y^f \;\in\; \R^{N_e \times N_e}.
\]
当 $R = \sigma_{\mathrm{obs}}^2 I_m$ 时，第 $(i,j)$ 元素为
\[
  (T_e)_{ij} = \delta_{ij} + \frac{1}{\sigma_{\mathrm{obs}}^2}
  \sum_{p=1}^{m} Y^f_{p,i}\,Y^f_{p,j}.
\]

\medskip\noindent
\textbf{步骤 4：求解分析均值权重。}
设创新向量 $\mathbf{d} = \y - \bar{\y}^f$。构造右端项
\[
  \mathbf{r} = (Y^f)\T R^{-1} \mathbf{d} \;\in\; \R^{N_e},
\]
求解 $N_e \times N_e$ 线性系统
\[
  T_e \w^a = \mathbf{r},
\]
得均值权重向量 $\w^a \in \R^{N_e}$。于是分析均值为
\[
  \bar{\x}^a = \bar{\x}^f + X^f \w^a.
\]

\medskip\noindent
\textbf{步骤 5：对称平方根分解。}
对 $T_e$ 做特征分解：
\[
  T_e = V \Lambda V\T,
\]
其中 $\Lambda = \operatorname{diag}(\lambda_1, \dots, \lambda_{N_e})$。
构造逆平方根矩阵
\[
  T = V \Lambda^{-1/2} V\T,
  \qquad
  \Lambda^{-1/2} = \operatorname{diag}(\lambda_1^{-1/2}, \dots,
                                         \lambda_{N_e}^{-1/2}).
\]

\medskip\noindent
\textbf{步骤 6：重构分析集合。}
第 $e$ 个分析成员由预测均值和扰动变换的组合给出：
\begin{equation}
  \x^{a,(e)}
  = \bar{\x}^f
  + X^f \w^a
  + \sqrt{N_e-1}\; X^f [T]_e,
  \label{eq:app-reconstruct}
\end{equation}
其中 $[T]_e$ 为 $T$ 的第 $e$ 列。利用 $A = \sqrt{N_e-1}\,X^f$，可将
\eqref{eq:app-reconstruct} 改写为直接在原始异常矩阵上操作的等价形式：
\[
  \x^{a,(e)}
  = \bar{\x}^f
  + \sum_{j=1}^{N_e} \left(
      [T]_{j e} + \frac{w_j^a}{\sqrt{N_e-1}}
    \right)
    \bigl(\x^{f,(j)} - \bar{\x}^f\bigr).
\]
此式同时更新 $h$、$hu$ 和 $hv$ 三个分量。

\medskip\noindent
\textbf{步骤 7：物理约束。}
与附录~\ref{app:enkf} 步骤 7 相同。

%---------------------------------------------------------------------
\bibliographystyle{plain}
\bibliography{ref}

\end{document}
